\paragraph{Teaching answer boundaries improves individual replies on pictured passages.}
A reader that selects a source interval answered both questions on 49 of the preceding 128 passages. The causal writer passed 15. We tested whether the reader could teach the writer where to copy and when to stop. Three copies of the mixed writer received the same 26,381 training examples. Ordinary replay supervised the complete answer. A second condition added the positions of matching source glyphs. A third divided that additional supervision equally between the annotated alignment and a frozen visual reader's probability distribution over intervals.

On 98 further Chinese passages with two different questions each, the parent, replay, alignment and distilled writers returned 60, 63, 68 and 72 exact answers out of 196 (Fig.~\ref{fig:interval-teacher}). The corresponding numbers of passages with both answers correct were 11, 11, eight and 11. Distillation gained four complete passages over replay and lost four others (paired bootstrap 95\% interval for the difference, $-6.1$ to $+6.1$ percentage points). The span reader passed 36. Erasing the question reduced the distilled writer's 11 complete passages to zero. Among its 124 incorrect replies with the question visible, 66 started with incorrect ink. Of the other 58, 20 extended an accepted answer and 15 stopped at a proper prefix.

On the earlier set of 359 human questions, distillation gave 144 exact answers, compared with 125 for replay and 242 for the reader. Distillation reproduced 192/192 historical requests in familiar wording and 56/192 with new wording. The same weights returned both orders of the supplied oracle and bronze forms on the known \textit{yan} card. On the two-paragraph date task with learned selection, it passed 20/48 twelve-question comparisons versus 23 for replay. On 571 separate DuReader development questions it exactly reproduced 10/177 references expressible as one contiguous span, zero of 94 requiring distributed ink and zero of 300 containing a content glyph absent from the input. An alternative wording can be meaningful despite missing the recorded reference. Distillation did not meet the prespecified complete-passage criterion for replacing the mixed writer.

In a separate batch-one comparison, Qwen3-8B Q8 answered 144/196 questions and passed 55/98 passages from the same sources supplied as text. Its mean warm response time was 263 ms. The distilled image writer took 46 ms per request on prepared images, or 56 ms including rendering. All 37 text-model layers ran on the same GPU. The two systems achieved different reading accuracy. The observed time ratio therefore measures serving latency at these operating points.

The same Chinese questions exposed a second limitation. Rerendering the pages and requests in bold reduced the distilled writer's score to 39/196 answers and seven complete passages. With Kai script, two passages contained unsupported font characters and were excluded. All four writers scored zero on the remaining 192 questions, even when the scorer accepted each answer glyph in either Kai or the training typeface.
